\section*{Chapitre \ref{ch:rs:cross-sectional-and-cross-asset-features} --- Variables transversales et variables multi-actifs}

\begin{solution}{exo:rs:cross-sectional-and-cross-asset-features:1}
Sans lui-même, la moyenne des pairs vaut $(1 + 2 + 3 + 4)/4 = 2.5\%$ et le \omterm{def:rs:cross-sectional-and-cross-asset-features:peer}{rendement relatif aux pairs} vaut 7.5~\%. Avec lui-même, la moyenne vaut 4~\% et
l'écart vaut 6~\%, soit $\frac45 \times 7.5~\%$, comme l'énonce la proposition pour $n = 5$.
\end{solution}

\begin{solution}{exo:rs:cross-sectional-and-cross-asset-features:2}
$10\,000 \times 133/486\,506 = 2.7$~; pour les 1\,000 premiers, 0.27.
\end{solution}

\begin{solution}{exo:rs:cross-sectional-and-cross-asset-features:3}
$0.1 \times (-6~\%) = -0.6~\%$ au total, soit $-0.6~\%/21 = -0.029~\%$ par jour (2.9 points de base) sur 21 jours.
\end{solution}

\begin{solution}{exo:rs:cross-sectional-and-cross-asset-features:4}
Avec $\sigma^2$ le taux de variance, le rendement de $B$ sur $((k-1)\Delta, k\Delta]$ est celui de $A$ sur $((k-1)\Delta - L, k\Delta - L]$. Cet intervalle
recouvre l'intervalle précédent de $A$, $((k-2)\Delta, (k-1)\Delta]$, sur une longueur $L$, et le même intervalle de $A$ sur $\Delta - L$. Les accroissements browniens
sur des intervalles disjoints sont indépendants, donc les covariances valent $\sigma^2 L$ et $\sigma^2(\Delta - L)$, et en divisant par la variance
$\sigma^2\Delta$ de chaque rendement donne $L/\Delta$ et $1 - L/\Delta$.
\end{solution}

\begin{solution}{exo:rs:cross-sectional-and-cross-asset-features:5}
Les rendements sur une grille plus fine que le temps entre transactions sont pour la plupart nuls, et leurs corrélations s'amenuisent (l'effet Epps du livre~4,
chapitre~21)~; une grille plus grossière que l'avance la fond dans le même intervalle. La somme de Hayashi--Yoshida utilise chaque rendement observé et
apparie ceux dont les intervalles se recouvrent, si bien que décaler une horloge de $\vartheta$ déplace les paires exactement, quel que soit l'espacement des transactions.
\end{solution}

\begin{solution}{exo:rs:cross-sectional-and-cross-asset-features:6}
Régresser la semaine du quintile des petites capitalisations sur sa propre semaine précédente et sur celle du quintile des grandes. Les deux régresseurs sont corrélés (les quintiles
évoluent ensemble au sein d'une semaine), si bien qu'une grande part de ce que la dernière semaine du quintile des grandes semble prédire est la dernière semaine du quintile des petites. Dans
la régression du chapitre, le coefficient du quintile des grandes n'a une statistique $t$ que de 1.9 de 1962 à 1987, de $-0.7$ de 1988 à 2006
et de 1.1 depuis.
\end{solution}

\begin{solution}{exo:rs:cross-sectional-and-cross-asset-features:7}
Avec une diffusion de 5 jours, l'IC de la variable mensuelle tombe à 0.008 ($t = 1.4$) alors que celui de la vérité terrain vaut 0.082~: à la fin du mois, la plus grande partie des
nouvelles du client a déjà atteint le fournisseur, et la variable regarde le mauvais mois. Avec 63 jours, l'IC de la variable vaut 0.031
($t = 4.8$) et celui de la vérité terrain 0.050 (la dérive plantée par mois est plus petite). La fenêtre d'une variable doit correspondre au temps que met l'information
à se diffuser, ce que mesure directement l'étude d'événement de la \cref{fig:rs:cross-sectional-and-cross-asset-features:event}.
\end{solution}

\begin{solution}{exo:rs:cross-sectional-and-cross-asset-features:8}
Tokyo clôture avant que New York n'ouvre~: le rendement du S\&P 500 à la même date n'est pas connu à la clôture de Tokyo, si bien que la corrélation est une
fuite d'horodatages (chapitre~3). La version exploitable régresse la prochaine séance de Tokyo sur la dernière séance achevée du S\&P, connue avant
l'ouverture de Tokyo, et doit battre le passé propre de Tokyo ainsi que les contrats à terme négociés la nuit.
\end{solution}

\begin{solution}{pb:rs:cross-sectional-and-cross-asset-features:1}
\begin{enumerate}
\item 0.65 fois par seconde dans la paire rapide, 0.16 dans la lente (environ toutes les six secondes).
\item Les prix milieux simulés évoluent par paliers et suivent une tendance sur quelques secondes, si bien que les rendements au niveau du tick sont autocorrélés et que leur somme des carrés
fausse la variance~; les rendements à la minute sont assez longs pour être presque non corrélés.
\item Pour 0.15 seconde~: la fonction reste plate à une seconde près environ de la latence, car chaque prix milieu suit son prix efficient avec une
dispersion de délais, et l'argmax tombe n'importe où sur le sommet plat.
\item 0.5 seconde. Elle suppose que les deux instruments réagissent à leurs prix efficients avec la même distribution de délais, ce qui rend la
corrélation croisée symétrique autour de la latence.
\item 1.17 pour 0.5 seconde~; 1.03 pour une latence nulle.
\item Le centre~: 0.053 seconde en moyenne, 0.2 au pire~; l'argmax~: 0.19 en moyenne, 0.5 au pire.
\item Leurs prix milieux changent toutes les six secondes~; un décalage d'une fraction de seconde modifie à peine les intervalles qui se recouvrent, si bien que la fonction est
plate sur plusieurs secondes, et son sommet n'est que du bruit.
\item Des instruments dont les prix changent plusieurs fois en 50 millisecondes, des horodatages issus d'une même horloge au point de réception, précis à une
petite fraction de ce délai, et assez de changements de prix pour que le sommet de la fonction se détache du bruit.
\item 0.59, 0.45 et 0.44.
\item Les prix milieux du simulateur atteignent un nouveau prix efficient par paliers sur quelques secondes, si bien que chaque rendement de prix milieu suit une tendance sur de courts intervalles~:
chaque série se prédit elle-même et l'autre.
\item La première ligne moins la seconde~: 0.13 à deux secondes, 0.17 à cinq, 0.10 à dix, 0.02 à trente.
\item À trente secondes, l'avance d'une demi-seconde ne représente qu'une petite part de chaque intervalle, et les tendances se sont épuisées à l'intérieur~; à 0.1
seconde, la plupart des intervalles ne contiennent aucun changement de prix.
\item Dans les deux cas, le passé propre du suiveur explique une grande partie de ce que le passé du meneur semble expliquer~; le test porte sur la contribution du meneur au-delà de celle-ci.
\item Probablement pas avec des ordres au marché~: une fourchette d'un tick est grande face à un mouvement d'une fraction de tick prédit avec une corrélation de 0.6,
et Huth et Abergel ont montré qu'une stratégie naïve aux ordres au marché ne pouvait pas tirer profit de l'avance à cause de la fourchette. La variable sert à
déplacer les cotations ou choisir le moment des ordres passifs.
\item IC de 0.035 avec $t = 5.6$~; celui de la vérité terrain vaut 0.074.
\item Le dernier mois du client mélange des nouvelles déjà parvenues au fournisseur et des nouvelles encore en diffusion, et ajoute les
mouvements de marché et de facteurs~; la vérité terrain est exactement la part encore à venir.
\item Quotidienne~: 3~; mensuelle~: 8~; au hasard~: 2.7.
\item \textbf{Résultat nommé.} Sur la paire rapide, le centre symétrique retrouve les latences plantées de 0 à 2 secondes à moins de 0.053 seconde en
moyenne (0.2 au pire)~; la prévisibilité que l'avance ajoute par rapport au sens inverse vaut 0.13 à deux secondes, 0.17 à cinq, 0.10 à dix
et rien à trente.
\item L'avance d'une demi-seconde~: la paire est connue, et une heure contient des milliers de changements de prix qui portent tous sur un seul paramètre. Le lien mensuel
se cache parmi 486\,506 paires candidates de 120 mois chacune~; sans la liste des clients, il n'est pas trouvé.
\item Elle prédit le suiveur au-delà de son propre passé, hors échantillon (out-of-sample), à une échelle de temps où elle survit au coût de négociation qu'elle implique,
avec un mécanisme qui explique pourquoi.
\end{enumerate}
\end{solution}

\begin{solution}{iq:rs:cross-sectional-and-cross-asset-features:1}
Par classification sectorielle, par description de l'activité, ou par les plus proches voisins en corrélation de rendements sur une fenêtre passée (telle que connue à
l'époque). Pièges~: inclure le titre dans sa propre moyenne de pairs (un facteur $(n-1)/n$ qui varie selon la taille des groupes), des groupes trop petits pour
avoir une moyenne stable, des pairs non cotés le jour concerné (\omterm{def:rs:point-in-time-data-and-the-biases:survivorship}{biais de survie}) et des classifications révisées après coup.
\end{solution}

\begin{solution}{iq:rs:cross-sectional-and-cross-asset-features:2}
Prendre les deux flux de cotations sur une même horloge au point de réception, calculer la covariance croisée de Hayashi--Yoshida des prix milieux avec l'horloge de l'ETF
décalée sur une grille de retards, et retenir le retard du pic, ou mieux le centre de symétrie de la fonction~; le vérifier sur
des données simulées à retard connu, par heure de la journée, et au regard du temps entre changements de prix, qui borne la résolution.
\end{solution}

\begin{solution}{iq:rs:cross-sectional-and-cross-asset-features:3}
Seulement si le passé des grandes valeurs prédit celui des petites au-delà de leur propre passé. Dans les quintiles de taille français, la corrélation hebdomadaire
du quintile des petites avec sa propre dernière semaine (0.37 de 1962 à 1987) était aussi forte qu'avec celle du quintile des grandes (0.31), et en contrôlant celle-ci
laisse au quintile des grandes une statistique $t$ de 1.9~; depuis 2007, la corrélation vaut 0.06.
\end{solution}

\begin{solution}{iq:rs:cross-sectional-and-cross-asset-features:4}
Surtout du bruit. Parmi 1\,000 titres, il existe 999\,000 paires ordonnées~; avec quelques années de données, la corrélation décalée d'un vrai lien n'est qu'une
fraction de l'écart type entre paires, si bien que les 100 premières sont les plus grands tirages de bruit. Sur le marché synthétique du chapitre, les
10\,000 premières des 486\,506 paires en contenaient 8 sur 133 liens plantés, contre 2.7 au hasard. Partir de \omterm{def:rs:cross-sectional-and-cross-asset-features:leadlag}{liens économiques}, et traiter la fouille
de données comme un problème de tests multiples.
\end{solution}

\begin{solution}{iq:rs:cross-sectional-and-cross-asset-features:5}
Les investisseurs à attention limitée suivent le titre, et non ses clients (Cohen et Frazzini)~; l'information sur les entreprises liées se diffuse
graduellement entre investisseurs spécialisés (Menzly et Ozbas). Tester avec une liste de clients en \omterm{def:rs:point-in-time-data-and-the-biases:lookahead}{données point-in-time}~: trier les fournisseurs sur le rendement de leurs clients
du mois passé, mesurer l'écart de rendement du mois suivant et sa statistique $t$, et vérifier la réponse en temps événementiel des fournisseurs
aux nouvelles des clients.
\end{solution}

\begin{solution}{iq:rs:cross-sectional-and-cross-asset-features:6}
Si le suiveur est le meneur retardé de $L$ et que les prix sont browniens, le rendement du suiveur sur un intervalle $\Delta \ge L$ a une covariance
$\sigma^2 L$ avec l'intervalle précédent du meneur et $\sigma^2(\Delta - L)$ avec le même~: la part prévisible vaut $L/\Delta$, et la
corrélation décalée vaut $L/\Delta$. Pour $\Delta < L$, toute la covariance se trouve dans des intervalles passés, répartie sur plusieurs. Le bruit et
l'autocorrélation propre du suiveur modifient les chiffres en pratique, et c'est pourquoi le chapitre mesure l'asymétrie entre les deux
sens.
\end{solution}
